% ============================================================
% Section 111: 顺磁杂质二能级系统比热（肖特基反常）
% ============================================================
\section{固体物理：顺磁杂质二能级系统比热}
\label{sec:schottky-anomaly}

\begin{modelbox}[题目：顺磁杂质对比热的贡献]
考虑一包含有顺磁杂质的晶体，杂质在给定的磁场中能级是 $\pm\varepsilon$。

(1) 写出磁场存在时，杂质对晶体配分函数的贡献；

(2) 写出杂质对比热贡献的相应表达式，并画出比热作为温度函数的曲线；

(3) 在低温极限，比较上述比热与声子比热（德拜模型）对温度的依赖关系，并物理解释；

(4) 在高温极限，做如 (3) 相同的比较。
\end{modelbox}

\begin{tipbox}[考点快速分析]
\begin{itemize}
    \item \textbf{二能级系统配分函数}：$Z_1=2\cosh(\varepsilon/k_BT)$
    \item \textbf{肖特基比热}：$c_V=k_B(\varepsilon/k_BT)^2\sech^2(\varepsilon/k_BT)$
    \item \textbf{低温声子比热}：德拜模型 $c_V\propto T^3$
    \item \textbf{高温声子比热}：杜隆--珀蒂极限 $c_V\to3Nk_B$
\end{itemize}
\end{tipbox}

\begin{derivationbox}[标准解题过程]
\textbf{(1) 杂质对配分函数的贡献}

每个顺磁杂质为二能级系统，能级 $-\varepsilon$ 和 $+\varepsilon$（$\varepsilon=\mu_BB$ 为塞曼劈裂）。单个杂质配分函数
\begin{equation}
Z_1=e^{-\varepsilon/k_BT}+e^{\varepsilon/k_BT}=2\cosh\!\left(\frac{\varepsilon}{k_BT}\right).
\end{equation}
设晶体中有 $N$ 个独立杂质，则总配分函数贡献 $Z=(Z_1)^N$。

\textbf{(2) 杂质对比热的贡献}

单个杂质的平均能量
\begin{equation}
\bar{E}=-\frac{\partial\ln Z_1}{\partial\beta}=-\varepsilon\tanh(\beta\varepsilon)=-\varepsilon\tanh\!\left(\frac{\varepsilon}{k_BT}\right).
\end{equation}

单个杂质对比热的贡献（$N$ 个杂质则乘以 $N$）：
\begin{align}
c_V&=\frac{d\bar{E}}{dT}
=\varepsilon\sech^2\!\left(\frac{\varepsilon}{k_BT}\right)\cdot\frac{\varepsilon}{k_BT^2}\\
&=k_B\left(\frac{\varepsilon}{k_BT}\right)^2\sech^2\!\left(\frac{\varepsilon}{k_BT}\right).
\end{align}
此即\textbf{肖特基比热}（Schottky anomaly）。

\textbf{曲线特征}：
\begin{itemize}
    \item 低温（$k_BT\ll\varepsilon$）和高温（$k_BT\gg\varepsilon$）下 $c_V$ 均趋于零；
    \item 在 $k_BT\sim\varepsilon$ 附近出现宽阔峰（肖特基反常），峰值约在 $k_BT\approx0.834\,\varepsilon$。
\end{itemize}

\textbf{(3) 低温极限的比较}

杂质比热：当 $k_BT\ll\varepsilon$ 时，$\sech(\varepsilon/k_BT)\approx2e^{-\varepsilon/k_BT}$，故
\begin{equation}
c_V\approx4k_B\left(\frac{\varepsilon}{k_BT}\right)^2e^{-2\varepsilon/k_BT}.
\end{equation}
比热以\textbf{指数形式}趋于零（激活型行为）。

声子比热（德拜模型）：当 $T\ll\Theta_D$ 时
\begin{equation}
c_V^{\text{ph}}\propto T^3,
\end{equation}
以\textbf{幂律形式}趋于零。

\textbf{物理解释}：杂质二能级系统存在能隙 $2\varepsilon$。低温下热激发能量远小于能隙，杂质几乎全部冻结在基态，需跨越能隙才能激发，比热呈指数衰减。声子系统是连续谱（无能隙），低温下仍有低频长波声子可被激发，激发声子数与 $T^3$ 成正比。

\textbf{(4) 高温极限的比较}

杂质比热：当 $k_BT\gg\varepsilon$ 时，$\sech(\varepsilon/k_BT)\approx1$，故
\begin{equation}
c_V\approx k_B\left(\frac{\varepsilon}{k_BT}\right)^2\propto\frac{1}{T^2}.
\end{equation}
比热随温度升高按\textbf{幂律衰减}。

声子比热（德拜模型）：当 $T\gg\Theta_D$ 时
\begin{equation}
c_V^{\text{ph}}\approx3Nk_B=\text{常数},
\end{equation}
趋于\textbf{杜隆--珀蒂极限}，与温度无关。

\textbf{物理解释}：杂质只有两个能级，高温下上下能级几乎等概率占据，系统接近饱和，能量随温度变化率减小，比热按 $1/T^2$ 下降。声子系统有无穷多能级，高温下每个模式均分得能量 $k_BT$，总内能线性增长，比热趋于常数。
\end{derivationbox}

\begin{formulabox}[答案汇总]
\begin{center}
\begin{tabular}{lll}
\toprule
\textbf{项目} & \textbf{杂质比热} & \textbf{声子比热（德拜）} \\
\midrule
表达式 & $k_B(\varepsilon/k_BT)^2\sech^2(\varepsilon/k_BT)$ & 德拜模型 \\
低温极限 & $\propto e^{-2\varepsilon/k_BT}\to0$（指数） & $\propto T^3\to0$（幂律） \\
高温极限 & $\propto1/T^2\to0$（幂律衰减） & $\to3Nk_B$（常数） \\
\bottomrule
\end{tabular}
\end{center}
\end{formulabox}

\begin{insightbox}[物理图像]
\begin{itemize}
    \item 二能级系统的肖特基比热呈现宽阔峰，是有限能级结构的典型热力学特征，广泛用于分析晶体场劈裂、塞曼劈裂等效应。
    \item 杂质比热与声子比热的本质区别源于能谱结构：杂质能谱有能隙且有限，声子能谱无能隙且无限。
    \item 实际晶体低温比热测量中，肖特基反常常叠加在 $T^3$ 声子背景上，通过分析可提取杂质能级劈裂信息。
\end{itemize}
\end{insightbox}
